% ============================================================================
% DreamingGoose: Staged Distillation from Autoregressive Transformers to
% Bidirectional Recurrent Diffusion Language Models
% ICML 2026 format, camera-ready style.
%
% Companion to main.tex ("Anatomy of Associative Recall in Fixed-State
% Recurrences"). That paper isolates mechanism on a toy bench; this one asks
% what survives when the same levers are applied to a real converted LM.
%
% Preliminary-results release (not a submission). Section 7 reports a run that
% is STILL EXECUTING and states no outcome for it.
%
% Build: pdflatex dreaminggoose && bibtex dreaminggoose && pdflatex dreaminggoose && pdflatex dreaminggoose
% ============================================================================
\documentclass{article}

\usepackage{microtype}
\usepackage{graphicx}
\usepackage{booktabs}
\usepackage[dvipsnames]{xcolor}
\usepackage{amsmath,amssymb}
\usepackage{mathtools}
\usepackage{enumitem}

\usepackage{hyperref}
\hypersetup{colorlinks=true,linkcolor=teal!60!black,citecolor=teal!60!black,urlcolor=teal!60!black}

\usepackage[accepted]{icml2026}
\makeatletter
\renewcommand{\Notice@String}{Preprint. Preliminary results, shared for discussion; not peer reviewed.}
\makeatother

% ---------- notation ----------
\newcommand{\gdn}{\texttt{gated\_deltanet}}
\newcommand{\disj}{\textsc{disjoint}}
\newcommand{\trained}{\textsc{trained-band}}
% ------------------------------

\icmltitlerunning{DreamingGoose: Staged Distillation to Recurrent Diffusion LMs}

\begin{document}

\twocolumn[
\icmltitle{DreamingGoose: Staged Distillation from Autoregressive Transformers\\
to Bidirectional Recurrent Diffusion Language Models}

\icmlsetsymbol{equal}{*}

\begin{icmlauthorlist}
\icmlauthor{Julian Boesch}{obit}
\icmlauthor{Andrew Wee}{obit}
\end{icmlauthorlist}

\icmlaffiliation{obit}{Obit Development}

\icmlcorrespondingauthor{Julian Boesch}{jboesch@obitmc.com}
\icmlcorrespondingauthor{Andrew Wee}{awee@obitmc.com}

\icmlkeywords{diffusion language models, linear attention, distillation, model conversion, curriculum learning, associative recall, Machine Learning, ICML}

\vskip 0.3in
]

\printAffiliationsAndNotice{}

\makeatletter
\gdef\@icmltitlerunning{DreamingGoose: Staged Distillation to Recurrent Diffusion LMs}
\makeatother

\begin{abstract}
Pretrained autoregressive Transformers represent enormous sunk compute. Whether that compute survives changing \emph{both} the architecture (attention $\to$ recurrent state) and the generative objective (next-token $\to$ iterative denoising) is open: existing work crosses one axis at a time. We convert attention teachers into attention-free bidirectional recurrent diffusion students in stages, and report what transfers. The headline is a negative that became a positive. In-context retrieval transfers at exactly $0.000$: a converted student inherits its teacher's knowledge but none of its ability to bind and retrieve within a sequence. A distance curriculum repairs this, but our first repaired model---$0.94$--$1.00$ retrieval, threefold its teacher---did not replicate, and its exact rerun returned $0.000$. The repair is a \emph{lock-in lottery}, and the fix is to close the loop: advancing the curriculum only while retrieval accuracy holds takes the same three seeds from one partial lock-in to $3/3$ at $0.94$--$0.99$, survives real text ($0.955$--$0.995$ on every cell), and carries unchanged to $8$B, where recall transfer rises from $0.000$ to a mean of $0.901$ and generalizes to twice the trained load and beyond the trained distance. At $8$B the residual failure is identified: bootstrap onset is heavy-tailed ($6.5$k, $11$k, ${>}16$k steps), so a fixed step budget truncates the tail rather than the method failing. One boundary resists everything we tried: every model that solves its trained token band reads chance on held-out bands, and an arm that randomizes bindings each batch---so identities cannot be memorized---still does, locating the gap in token \emph{coverage} rather than memorization. We report the withdrawn number, the three pool-arm failure modes, and an in-flight block-diffusion conversion for which we state no result, because the discipline this pipeline most needs is distinguishing a lucky draw from a method.
\end{abstract}

% ============================================================================
\section{Introduction}\label{sec:intro}

Two conversion literatures exist and they do not intersect. One adapts pretrained autoregressive Transformers into \emph{diffusion} language models while keeping the Transformer backbone \citep{gong2025diffullama,ye2025dream}. The other distills attention into \emph{linear-time recurrence} while keeping next-token prediction \citep{goldstein2025radlads}. Crossing both axes at once---attention to recurrent state \emph{and} left-to-right to full-sequence denoising---is the regime an attention-free diffusion LM actually requires, and it is the regime nobody has reported.

We do it in stages, which is the paper's method as much as its engineering: distillation of the attention teacher into a recurrent student, bidirectional expansion, then diffusion continual pretraining. Staging exists so that capability loss can be \emph{attributed}. When something fails to transfer, we can say which stage lost it.

The finding that organizes everything below is a clean dissociation. Converted students keep what the teacher knows and lose what the teacher can \emph{do within a sequence}: on a teacher-gated multi-query associative recall probe, where the attention teacher scores $0.335$--$0.350$, the converted student scores $0.000$---not degraded, absent. This is the practical face of the interference wall analyzed on a toy bench in our companion paper \citep{anatomyrecall}, and the repair that works there---a table-to-query distance curriculum---is what we port here.

Porting it taught us the lesson this paper is really about. The first converted model to which we applied the curriculum reached $0.94$--$1.00$ retrieval across pair loads and needle distances, roughly threefold its own attention teacher. That number does not replicate. An exact-configuration rerun returned $0.000$; the run predated the pipeline's seed flag, so its draw cannot be reproduced by seed; and its checkpoint was lost to a storage-quota fault. We withdraw it as a result and keep it as one cell of a rate. What replaced it is better: the failure mode is legible, and closing the curriculum's control loop removes it.

\textbf{Contributions.}
\begin{enumerate}[leftmargin=1.4em,itemsep=2pt]
\item \textbf{A staged conversion pipeline that crosses both axes} (\S\ref{sec:pipeline}), with a stage-attributable account of what transfers: teacher knowledge does, in-context retrieval does not ($0.000$, at both $1.7$B and $8$B).
\item \textbf{The conversion-scale lottery, and its closed-loop fix} (\S\ref{sec:lottery}). Under an open-loop distance ramp the repair lands on one seed in three. A \emph{success-gated} ramp---the gap cap advances only while an EMA of retrieval accuracy clears a threshold---lifts the same three seeds to $3/3$ at $0.94$--$0.99$, and holds on real text ($0.955$--$0.995$ on every evaluated cell).
\item \textbf{The recipe carries to 8B, and the residual failure is named} (\S\ref{sec:8b}). Recall transfer $0.000\to0.901$ mean, generalizing to $N\le32$ and gaps to $224$; $2/3$ seeds lock in. Every locked-in run shows slow bootstrap then avalanche, with onset at $6.5$k, $11$k and ${>}16$k steps---heavy-tailed, so fixed budgets truncate the tail. The indicated lever is a success-gated \emph{budget}, not more steps.
\item \textbf{A boundary that survives every intervention we tried} (\S\ref{sec:coverage}). Held-out token bands read chance for every solved arm at both scales. A role-randomized arm excludes memorized bindings as the cause, locating the gap in token coverage; three curriculum policies over randomized pools fail by three distinct mechanisms (lockout, learn-then-collapse, non-bootstrap).
\item \textbf{Negative and in-flight results reported as such} (\S\ref{sec:blockdiff}, \S\ref{sec:limits}): a representation-alignment auxiliary that we armed and then disarmed on evidence, and a block-diffusion conversion still executing, for which we state no outcome.
\end{enumerate}

% ============================================================================
\section{Related Work}\label{sec:related}

\textbf{AR-to-diffusion adaptation.} \citet{gong2025diffullama} and \citet{ye2025dream} adapt pretrained autoregressive Transformers into masked diffusion LMs, establishing that the objective change is survivable and far cheaper than pretraining. Both retain attention. \textbf{Attention-to-recurrence distillation.} RADLADS \citep{goldstein2025radlads} converts attention into RWKV-7-style recurrence at multi-billion scale on a few hundred million tokens, retaining next-token prediction. \textbf{Recurrent denoisers.} DiffuMamba \citep{singh2025diffumamba} shows scan backbones are viable denoisers; the design space beyond Mamba, and specifically the rank-1 state update of \citet{peng2025rwkv7} and \citet{yang2024gateddeltanet}, motivates the student architecture we convert into. \textbf{Block diffusion} \citep{arriola2025bd3lm} interpolates between the autoregressive and diffusion objectives, and is the objective of the conversion reported as in-flight in \S\ref{sec:blockdiff}. \textbf{Recall in recurrence.} That fixed-state recurrences underperform attention on associative recall is well documented \citep{arora2023zoology,arora2024based}; that the deficit is substantially a \emph{learnability} problem rather than a capacity limit is the reading our companion paper argues on a toy bench \citep{anatomyrecall} and that this paper tests at scale, alongside convergent evidence \citep{okpekpe2025recall,blouir2024birdie}.

% ============================================================================
\section{The Conversion Pipeline}\label{sec:pipeline}

The pipeline takes an attention teacher to an attention-free bidirectional recurrent diffusion student in three stages, each with its own objective and its own frozen set.

\textbf{Stage 1 --- architectural transfer.} Attention layers are replaced by gated delta-rule recurrent cells and trained to match the teacher's hidden states and logits, following RADLADS \citep{goldstein2025radlads}. The objective is still next-token prediction; only the architecture moves.

\textbf{Stage 2 --- bidirectional expansion.} A backward recurrent stream is introduced, initialized from mirrored forward weights, and fused with the forward stream by a merge module conditioned on the diffusion timestep. The forward stream is frozen here, so this stage learns fusion and nothing else.

\textbf{Stage 3 --- objective transfer.} Continual pretraining under a masked discrete diffusion objective \citep{sahoo2024mdlm,austin2021d3pm}, optionally mixed with the retrieval curriculum described below.

Staging is what makes the central dissociation attributable. We evaluate retention with a teacher-gated multi-query associative recall probe: key--value pairs are laid out in a table, a query is issued after a distractor haystack, and the model must return the bound value. \emph{Gated} means the teacher must itself pass before a student number is admitted, which excludes probe artifacts as an explanation for student failure---the same discipline our companion paper applies to its state-tracking guardrail \citep{anatomyrecall}.

The result of that evaluation is stark and is the reason for everything that follows. The attention teacher scores $0.335$--$0.350$. The converted student, at both $1.7$B and $8$B, scores $0.000$ against a marginal of ${\approx}0.023$: not degradation, but absence. Perplexity-shaped capability transfers; the in-sequence binding faculty does not. Under stage-3 continual pretraining alone it does not return.

\textbf{The curriculum, and one prediction it makes.} The repair ported from the toy bench samples the table-to-query gap rather than fixing it at maximum, so the model first solves the adjacent case and extends outward. Injected into stage 3, it comes with a falsifiable prediction from the mechanism account: if the circuit is \emph{mark-and-route}---mark the binding on the way past, route it to the query---then the routing direction must be trainable, so freezing the forward cells should teach the answer format while never acquiring bindings. That is what happens. With forward cells frozen (the pipeline default) the model learns to emit well-formed answers and never binds; unfreezing them is what admits the repair. We report this contrast as the one part of the original at-scale result that replicated cleanly.

% ============================================================================
\section{The Lottery, and the Closed-Loop Fix}\label{sec:lottery}

\begin{table}[t]
\centering
\caption{\textbf{Conversion-scale replication.} Retrieval on the trained token band, per seed (mean over $N\in\{4,8,16\}\times\text{gap}\in\{0,64,160\}$; marginal ${\approx}0.023$). \disj{} is the identical evaluation with key/value symbols shifted out of the trained band. The open-loop row's single non-zero seed is a \emph{partial} lock-in ($0.40$--$0.70$ across cells); closing the loop lifts all three of those same seeds to ceiling. Every row---including every solved one---reads $0.000$ disjoint.}
\label{tab:atscale}
\vskip 0.1in
\footnotesize
\setlength{\tabcolsep}{4pt}
\begin{tabular}{@{}llcc@{}}
\toprule
setting & lock-in & per-seed (main) & disj. \\
\midrule
1.7B, open-loop     & 1/3          & $0.00$/$0.00$/$0.55$ & $0.000$ \\
1.7B, \textbf{gated} & \textbf{3/3} & $\mathbf{0.97}$/$\mathbf{0.98}$/$\mathbf{0.97}$ & $0.000$ \\
1.7B, gated, wiki   & 1/1          & $0.955$--$0.995$     & $0.000$ \\
8B, gated           & 2/3          & $0.90$/$0.78$/chance & $0.000$ \\
\bottomrule
\end{tabular}
\end{table}

\textbf{The withdrawn number.} Our first curriculum-repaired conversion reached $0.94$--$1.00$ across $N\in\{4,8,16\}$ and gaps to $160$ tokens of real text, with shuffled-binding controls collapsing and a modest denoising cost ($-0.06$ against a stock control). Read as a result, it says the toy repair transfers at full strength. It does not. Three follow-ups seeded from the same configuration drew $0.000$, $0.000$, and a partial $0.40$--$0.70$; an exact-configuration rerun of the original drew $0.000$. The original also predates the pipeline's seed flag, so its draw is unreproducible by seed, and a storage-quota fault destroyed its checkpoint before the discrepancy was noticed. We state this plainly because the same lesson arrived independently on the toy bench, where a two-seed directional margin reversed under ten-seed replication \citep{anatomyrecall}: in this training regime, low-seed means are not evidence.

\textbf{What the failure actually is.} The open-loop curriculum advances the gap cap on a schedule---a clock. A run that has not yet learned to bind at short gaps is dragged to long ones anyway, and trains the rest of its budget at distances it cannot solve. Whether a seed escapes is close to a coin flip, which is precisely the lock-in lottery structure the toy suite reports.

\textbf{Closing the loop.} We make the schedule conditional: the gap cap advances only while an exponential moving average of retrieval accuracy stays above a threshold ($0.6$). No run is dragged up the ramp while failing; a run that stalls simply trains longer where it is. On the same three seeds that drew $\{0.000, 0.000, 0.40\text{--}0.70\}$ open-loop, the gated ramp locks in $3/3$ at $0.94$--$0.99$ (per-seed means $0.971$, $0.977$, $0.974$), distance-robust at the maximum evaluated gap, with shuffled-binding controls collapsing (Table~\ref{tab:atscale}). Telemetry confirms the mechanism rather than merely the outcome: all three arms climbed with the EMA at $0.98$--$0.99$ throughout. Final caps ($0.66$--$0.87$) never reached $1.0$, yet evaluation is already ceiling-robust at maximum gap---the tail of the ramp is unnecessary at this sequence length.

\textbf{It survives real text.} Every earlier attempt on a real corpus had read chance, leaving open whether the repair was an artifact of synthetic haystacks. Under the gated ramp a wikitext arm reaches $0.955$--$0.995$ on \emph{every} $N\times\text{gap}$ cell, including $N{=}16$ at gap $160$ ($0.975$), with clean controls. Its telemetry explains the earlier failures: bootstrap is merely \emph{slower} on real text (EMA $0.04$ at $3$k steps, avalanche by $5.4$k, clean ramp to cap $0.91$). Real text delays lock-in; it does not prevent it. The conversion-scale recipe is therefore curriculum $+$ trainable forward cells $+$ success-gated ramp, and the claim it supports is a lock-in \emph{rate}.

% ============================================================================
\section{Scaling to 8B}\label{sec:8b}

Applied unchanged to an $8$B gated-delta student, the recipe lifts recall transfer from $0.000$ to $0.69$--$0.99$ (mean $0.901$) on the same grid, with a bidirectional denoising gain of $+0.117$ over the causal stream---the query-first read that a bidirectional denoiser gets structurally, cashed out at scale. A wide evaluation of that model holds at twice the trained pair load and past the trained distance ($N\le32$, gap $\le224$: $0.49$--$0.99$, mean $0.786$), so both capacity and distance generalize beyond the training distribution rather than the model having memorized its evaluation grid.

\textbf{The residual lottery has a name.} Across three $8$B seeds the tally is $2/3$. The miss is not mysterious: every locked-in run, at both scales, shows the same profile---a long dead phase, then an avalanche in which retrieval NLL collapses within roughly a thousand steps, then a clean gated ramp. What varies is \emph{when}. Observed onsets are $6.5$k, $11$k and ${>}16$k steps, with a randomized-pool arm at $13$k and the wikitext arm at $5$k. The distribution is heavy-tailed, and every run here had a fixed step budget, so the budget truncates the tail. The consequence is visible in the amplitudes: seed 1 avalanched at $11$k and finished with a weaker floor (worst cell $0.35$) than seed 0, which avalanched at $6.5$k (worst cell $0.69$)---later onset means less post-avalanche training. The indicated lever is therefore a success-gated \emph{budget}---train until avalanche, then a fixed number of post-avalanche steps---rather than a larger fixed total. It is designed and not yet run; we report $2/3$ as it stands.

\textbf{A second architecture, stable but slower to bootstrap.} An RWKV-7 student at $8$B trains stably under the same recipe and shows a positive bidirectional denoising gain ($+0.053$, against $-0.024$ in an earlier unstable configuration), but does not avalanche within $16$k steps. Given the onset distribution above we read this as underpowered rather than negative: RWKV's retrieval bootstrap is slower than the gated delta-rule cell's everywhere we have measured, and it needs the gated budget before a verdict is available.

% ============================================================================
\section{What Does Not Transfer: Token Coverage}\label{sec:coverage}

Every arm that solves the trained-band evaluation---at $1.7$B and $8$B, including all three gated-ramp arms and the full-strength wikitext arm---scores $0.000$ on a \disj{} evaluation that shifts the key/value symbols deeper into the tokenizer and changes nothing else. What the curriculum installs is a binding circuit over the symbol pool that appeared in retrieval episodes, not an abstract retrieval faculty.

The obvious explanation is memorization of fixed bindings, and it is wrong. An arm trained with key/value roles randomized every batch---so no fixed identity \emph{can} be memorized---still bootstraps (late, at ${\sim}13$k steps, with the same avalanche profile) and still reads $0.000$ disjoint. It also generalizes distance far past its own trained cap: the cap reached only ${\approx}40$ tokens, yet gap-$64$ performance matches gap-$0$ and gap-$160$ sits well above chance. So the circuit is not memorizing pairs, and it is not distance-bound. The gap is which \emph{embeddings} ever participated in a retrieval episode.

\textbf{Three failure modes, one root cause.} Randomizing the pool interacts badly with the gate, and mapping that interaction was necessary to avoid mistaking it for a refutation. At gate $0.6$ the pool task starts below threshold, the cap never advances, and the arm trains only at gap ${\approx}0$ (\emph{lockout}). At gate $0.25$ the cap advances on a weak signal and then the EMA collapses; because the cap is a ratchet with no retreat, all subsequent training runs at impossible gaps and the circuit never recovers (\emph{learn-then-collapse}). With a backoff thermostat that lets the cap retreat, a pool of $512$ never bootstraps at all within budget, while a pool of $128$ does---at $13$k steps (\emph{non-bootstrap}, and plausibly just slower still). Three curriculum policies, three mechanisms, one cause: the randomized-pool task has a much later bootstrap onset, and every policy we tried was budget-limited rather than policy-limited.

\textbf{Status of the fix.} The diagnosis points at pool-size annealing---ride the avalanche on a narrow pool, then widen coverage. At $8$B, the one annealing run never bootstrapped its first tier inside a $24$k-step budget, so the tiers never engaged. Annealing is therefore \emph{untested, not refuted}, and it is blocked on the same success-gated budget that \S\ref{sec:8b} identifies. We attach this boundary to every at-scale number we report.

\textbf{A hybrid escape hatch.} Because token coverage resisted every training-side intervention we tried, it is worth naming as a reason to keep some attention. An attention-retaining conversion (keeping a sparse set of attention layers) now passes end-to-end validation---initialization agreement $0.490$ refined to $0.599$, perplexity retention $0.703$---but we report it as a working pipeline, not as a measured retrieval result.

% ============================================================================
\section{In Flight: Block-Diffusion Conversion at 8B}\label{sec:blockdiff}

The conversions above use full-sequence masked diffusion. A block-diffusion objective \citep{arriola2025bd3lm}---which interpolates toward autoregression by denoising within blocks conditioned on committed prefixes---is a better match for an instruction-following target, and we are converting a Qwen2.5-Coder $8$B teacher into a $3{:}1$ gated-delta\,:\,attention hybrid under it. \textbf{We state no result for this run.} At the time of writing it is mid-pretraining, and the honest report is the design, two negatives it has already produced, and the fact that its evaluation is not yet in hand.

\textbf{Construction.} Each step builds clean and noised views of the same batch. The clean stream is an ordinary causal pass of the same network; the noised stream conditions, per block, on the clean stream's state for strictly-previous blocks and on itself within-block. Setting the diffusion block equal to the recurrent chunk makes this nearly free: the kernel already materializes each chunk's incoming state as a gradient residual, so promoting it to an output yields exactly ``clean recurrent state at the start of block $n$'' at zero forward cost, and the noised stream needs no scan at all.

\textbf{Negative 1: a representation-alignment auxiliary, armed then disarmed.} We added a term aligning the student's clean-stream hidden states to the frozen teacher's at layer boundaries, on the theory that it would preserve teacher structure through the objective change. It was armed by default and is now off. The measured reason: the cosine term was already ${\approx}94\%$ saturated at initialization, so it had almost no gradient to give; autoregressive retention did not improve; and it inflated logit scale roughly twofold. We report it because an auxiliary that looks well-motivated, costs ${\approx}10$--$15\%$ step time, and does nothing is exactly the kind of thing that survives in a pipeline unexamined.

\textbf{Negative 2: learning rate in the mature regime.} The run diverged twice at learning rate $3\mathrm{e}{-4}$, both times deep into training (around steps $50$--$52$k and $56$k) and both times with an identical signature: escalating oscillation in accepted loss, spike-recover-collapse, accuracy to ${\approx}0.02$. Notably every individual step stayed under the loss-spike guard's threshold, so a step-level guard registered nothing wrong---liveness is not health. The cure was the same one the $1.5$B tail needed: $2\mathrm{e}{-4}$. We record this because the failure is invisible to the obvious monitor.

% ============================================================================
\section{Limitations}\label{sec:limits}

\begin{enumerate}[leftmargin=1.4em,itemsep=2pt]
\item \textbf{Token generality is unsolved.} Every solved arm reads $0.000$ on held-out token bands (\S\ref{sec:coverage}). We make no claim that the curriculum installs a general retrieval faculty, and the annealing fix that our own diagnosis motivates is untested rather than refuted.
\item \textbf{Three seeds, not ten.} Conversion-scale rows rest on three seeds per setting (one, for the wikitext arm). Given that this paper's central correction is about low-seed evidence, we hold our own rates to the same standard: $3/3$, $2/3$, and $1/1$ are small-$n$ rates, and no pairwise contrast here is significance-tested.
\item \textbf{One withdrawn number.} The first at-scale result is withdrawn (\S\ref{sec:lottery}) rather than silently restated. Its checkpoint is unrecoverable, so it can never be re-examined---an argument for artifact discipline, not a scientific defence.
\item \textbf{Fixed budgets confound several negatives.} The $8$B seed-2 miss, the RWKV-7 non-bootstrap, and the entire pool-annealing arm are all consistent with budget truncation of a heavy-tailed onset rather than method failure. Until the success-gated budget is run, those three are open, not settled.
\item \textbf{One teacher family per scale, and one probe family.} Retrieval is measured with a single multi-query associative recall probe in a curriculum-adjacent format. Held-out bindings are used, but the format proximity between training and evaluation is real and we do not claim broad downstream retrieval.
\item \textbf{The block-diffusion conversion has no result} (\S\ref{sec:blockdiff}). It is reported for its design and its two negatives only.
\end{enumerate}

% ============================================================================
\section{Discussion}\label{sec:discussion}

\textbf{What conversion actually costs.} The staged pipeline gives a clean answer to the question that motivates it. Teacher knowledge survives crossing both axes; the in-sequence binding faculty does not survive either. It is not degraded by conversion---it is absent, at $0.000$, and no amount of continual pretraining under the diffusion objective brings it back. It has to be retrained deliberately, and the mechanism story from the toy bench \citep{anatomyrecall} says why: the wall is a training-coverage gap, so only training that covers the gap can close it.

\textbf{Closed-loop curricula are the transferable lesson.} The single most useful thing we learned is not the recall number but the schedule. An open-loop curriculum is a clock; a closed-loop one is a controller, and the difference is $1/3$ versus $3/3$ at $1.7$B on identical seeds. The same distinction predicts the toy bench's length boundary, where a time-based ramp collapses to $0/6$ at $L{=}512$ \citep{anatomyrecall}. Wherever a curriculum has a difficulty knob and a measurable competence signal, gating the former on the latter appears to convert a lottery into a method. The natural extension---and the one all our open negatives point at---is to gate the \emph{budget} the same way, since a fixed step count is itself an open-loop decision about a heavy-tailed random variable.

\textbf{What the token-coverage boundary means.} A converted model whose retrieval was trained over a narrow symbol pool reads chance outside it, and this survived pool randomization, three gate policies, and a scale change. Taken at face value it says curriculum-installed retrieval is a circuit over participating embeddings, and that coverage is a first-class axis of the training distribution rather than an afterthought. It also supplies an unanticipated third reason to keep attention in a hybrid---alongside capacity and length---which is a more interesting design conclusion than the one we set out to reach.

\section*{Reproducibility Statement}
Per-arm telemetry (gate cap, retrieval EMA, avalanche onset), evaluation JSONs for every arm including the failed ones, and checkpoints for all arms except the withdrawn original are retained in object storage. The withdrawn run's checkpoint was destroyed by a storage-quota fault before the replication discrepancy surfaced; the incident is the reason every later arm is pushed and its remote listing verified before the run advances.

\section*{Impact Statement}
This work concerns making pretrained language models cheaper to re-architect, which lowers the compute cost of deploying capable models and correspondingly lowers the barrier to deploying them without scrutiny. The retrieval boundary we report is also a safety-relevant one: a converted model can look competent on its training distribution and read at chance just outside it, which is a failure mode that ordinary perplexity-shaped evaluation does not surface.

\bibliographystyle{icml2026}
\bibliography{references}

\end{document}
